\chapter[Aus Hamburg, Mecklenburg und …]{Aus Hamburg, Mecklenburg und Schleswig-Holstein}
\quellseite{5}

\section{I.}

„Das Himmelreich ist gleich einem Senfkorn, das ein Mensch nahm und säete es auf seinen Acker. Welches das kleinste ist unter allen Samen; wenn es aber erwächst, so ist es das größte unter dem Kohl, und wird ein Baum, daß die Vögel unter dem Himmel kommen und wohnen unter seinen Zweigen“, Matth. 13, 31. f. Als der HErr JEsus einst vor nun fast 1900 Jahren dies Gleichniß dem Volke Israel vorlegte, war es ein wahres Wort; denn es kam aus dem Munde der ewigen Wahrheit. Aber die Kirche aller Zeiten hat dies Wort ihres HErrn und Meisters auch herrlich bestätigt gefunden. Die Kirche des alten Bundes, die Synagoge, aus welch bescheidenen Anfängen ist sie herausgewachsen! Da waren die Patriarchen, die da predigten von des HErrn Namen; da war Moses, der ihnen das Gesetz, das Wort des Mundes Gottes, gab und den Samen der himmlischen Wahrheit ins Herz pflanzte. Und wie gering und unscheinbar war die Schaar der Jünger des HErrn, die den Grundstock der neutestamentlichen Kirche bildeten; aber wir sehen sie nach der Ausgießung des Heiligen Geistes als Zeugen JEsu zu Jerusalem und in ganz Judäa und Samaria und bis an das Ende der Erde, und überall, wohin sie ihren Fuß setzten, blühte die Kirche JEsu Christi, des auferstandenen Siegesfürsten, mächtig empor. Wie armselig und verachtet war jener Mönch in Wittenberg, der das unter dem Schutt römischer Irrlehre vergrabene Bibelbuch wieder hervorholte und zu Ehren brachte. Aber Luthers Hammerschläge an die Schloßkirche, begleitet von der Stimme des Engels: „Fürchtet Gott und gebet ihm die Ehre!“ hallten wieder in der ganzen Welt, in den Herzen Tausender und Abertausender von Kindern Gottes. Die Missouri-Synode, die wie ein Wunderbau vor unsern Augen dasteht, begann so ärmlich und dürftig, daß die Feinde wohl lächelnd fragen mochten: „Was, meinst du, will aus dem Kindlein werden?“ Die evangelisch-lutherische Freikirche Deutschlands, der niemand im Hinblick auf
\quellseite{6}
den Druck, unter dem sie steht, ein langes Leben prophezeit hätte, die noch tagtäglich eine solch ärmliche Entwickelung vor Augen hat, sie ist im Laufe von 25 Jahren doch zu einer Gemeinschaft erstarkt, die im reichen Segen wirkt. Ja, das Reich Gottes beginnt überall klein, es ist ein Senfkorn, das durch das Evangelium in der Menschen Herzen gesenkt wird und durch den Thau und Regen vom Himmel wächst und zunimmt, so daß endlich die Zweige dieses göttlichen Reiches sich über die ganze Erde erstrecken.

So hat auch die Aussaat, welche in Hamburg vor wenigen Jahren begonnen wurde, die Hoffnung nicht betrogen. Zwar ist unsere Hoffnung nicht in der Weise, wie wir wünschten und baten, in Erfüllung gegangen, daß nämlich der treulutherischen Gemeinde in dieser Hafenstadt am Elbstrom ein so rasches Emporblühen beschieden sein möchte, daß sie nach kurzer Zeit schon auf eigenen Füßen stehen könnte; nein, Gott ist auch hier seine eigenen Wege mit uns gegangen. Durch die Hamburger Mission hat unsere Kirche Eingang gefunden in Schleswig-Holstein und Mecklenburg, und unsere Gemeinde hat sich genöthigt gesehen, ein eigenes Pfarramt in Flensburg aufzurichten, und bittet die lieben Glaubensbrüder um Unterstützung für dieses Werk. Wenn wir an eine Theilung der Hamburger Parochie herantraten, so geschah dies nicht, um meine Person zu entlasten, obwohl die an mich gestellten Anforderungen auf dem weitausgedehnten Felde hier und da Strapazen mit sich brachten, die meiner Gesundheit nicht gerade förderlich waren. Ich weiß, meine Amtsbrüder liegen auch nicht auf Rosen gebettet, und wenn ich erst an die Reiseprediger in den Prairieen Americas denke oder an einen australischen Studiengenossen, der mir von seinen Wagentouren erzählt, so freue ich mich doch, zumeist wenigstens in einem deutschen Eisenbahnwagen dritter Klasse über den Erdboden befördert zu werden. Auch ist mir jeder einzelne meiner Missionsposten so ans Herz gewachsen, daß es mir sehr schwer wurde, die Arbeit hier oder da aufgeben zu müssen. Aber das Wohl der Gemeinden erforderte eine Theilung. Wollte ich ihren Bedürfnissen einigermaßen genügen, so mußte ich den Charakter eines Reisepredigers im vollen Sinne des Wortes annehmen. Hiergegen sprachen jedoch verschiedene Gründe: Es würde damit Hamburg der Vorzug, den es bei Aufrichtung eines selbständigen Predigtamtes überkommen hatte, entzogen worden sein, dies aber zum Schaden der Mission. Wohl hat auch die Hamburger Gemeinde die segensreiche Einrichtung der Lesegottesdienste und hat sich immer fleißig und gern zu den Lesegottesdiensten eingefunden. Doch wir dürfen nicht vergessen, daß die dortige Gemeinde eine Missionsgemeinde ist. In Hamburg sowohl wie auf dem ganzen übrigen Felde ist Missionsarbeit zu verrichten; überall muß ich neben der Sorge, die mir von Gott anvertrauten Seelen zu erbauen, auch darauf bedacht sein, Außenstehende unter den Schall des Wortes zu bringen. Und ist nun auch das vom Lesepulte vernommene Wort ebenso lebendig und kräftig wie dasjenige, welches von der Kanzel erschallt, kann es daher auch als eine lebenschaffende Kraft Gottes selig machen alle, die daran glauben, so weiß doch jeder Missionar, wie wenig die Bitte, die Lesegottesdienste zu besuchen, Beachtung findet Seitens solcher, die noch fern stehen. Allein nicht nur die Rücksicht auf unsere Mission erheischt mehr Predigtgottesdienste, nein, auch unsere eigenen Gemeindeglieder bedürfen ihrer zum Wachsthum in der Erkenntniß und in der Heiligung, zur Stärkung und Erhaltung im Glauben. Sie sehen sich auf Schritt und Tritt von Gefahren und Versuchungen umringt und bedroht, so daß sich in ihren Herzen wohl die Sehnsucht regen mag nach dem Wort, das Gott in dem Amt, das die Versöhnung predigt, seinen Kindern beschert hat.

Ferner zeigte es sich immer klarer, daß durch das Fehlen der dänischen Sprache in unsern Gottesdiensten in Nord-Schleswig der Strom des Evangeliums erheblich eingedämmt wird. Nicht nur würden dortige Einwohner dänischer Zunge sich bewegen lassen, Gottesdiensten in ihrer Sprache beizuwohnen, sondern es ist auch von solchen Außenstehenden der Wunsch geäußert, man möchten es ihnen doch ermöglichen, an der Predigt des reinen Wortes Theil zu haben. Nun nahm aber meine Thätigkeit in Hamburg so sehr alle meine Kräfte in Anspruch, daß ich immer wieder vom Studium der dänischen Sprache abgezogen wurde und erkannte, nur durch ständigen Aufenthalt im Norden hierin Wandel schaffen zu können.

Auf dem Felde unserer Hamburger Gemeinde begegnen wir allenthalben einem zwar langsamen, aber stetigen Wachsthum. Die Seelenzahl der Gemeinde hat das erste Hundert um 30 Seelen überschritten. In Rostock ist durch Lehrer Güllhoffs Eintritt in die Auswanderermission zu Bremen unserer dortigen Wirksamkeit scheinbar ein Ziel gesetzt, aber, Gott Lob, nur scheinbar. Der Sauerteig des Wortes säuert fort. Schon zeigt sich wieder Frucht. Dr. Walthers Schriften werden hier und da gelesen und erquicken die Seelen. Unser Zeugniß ist bis in die Spitze der mecklenburgischen Aristokratie gedrungen. Wohl ist hier die Zukunft dunkel, aber ich lege sie vertrauensvoll in die Hände dessen, der Menschenherzen wenden kann, wie es ihm wohlgefällt.

Die Gemeinde in Hamburg hat im verflossenen Jahre (1902) einen ungeahnten Aufschwung genommen. Die anfängliche Zeit der Dürre ist vorübergegangen, und unter dem erfrischenden Thau und Regen des Himmels und unter den belebenden Strahlen der himmlischen Sonne wächst und grünt und blüht es mächtig in der Gemeinde. Zwanzig Seelen durften wir aufnehmen, und als ich am letzten Trinitatisfeste wieder einen Gottesdienst in Hamburg hielt, wurde eine Wittwe mit zwei Kindern der Gemeinde hinzugefügt. Durch diesen Zuwachs ist die Finanzkraft der Gemeinde so vermehrt worden, daß Hamburg im Stande ist, 600 Mark und darüber aufzubringen; und am Ende des vorigen Rechnungsjahres hatten wir uns gefreut zu hören, daß die Gemeindeschuld auf 55 Mark zusammengeschmolzen war. So hat Gott das treue, mit schweren Geldopfern verbundene Ausharren der sächsischen Freikirche auf diesem wichtigen Posten reich gesegnet, und es ist der Hamburger Gemeinde jetzt eine Aussicht eröffnet, die zu den schönsten Hoffnungen berechtigt. Eine Störung für die Hamburger Mission ist natürlich der gerade in das Wachsthum schlagende Personenwechsel. Allein die Erkenntniß, daß unter den schwierigen Verhältnissen in Schleswig-Holstein ein Mann vonnöthen sei, der sich daselbst bereits einigermaßen eingearbeitet habe, ließ die lieben Hamburger einwilligen, daß ich nach Flensburg übersiedelte. Und Gott, der das Wort gibt mit großen Schaaren Evangelisten, hat ja auch die lieben Hamburger nicht Waisen gelassen, er hat ihnen in P. Uplegger einen tüchtigen und treuen Seelsorger geschenkt und wird auch weiter zu ihnen kommen und sie segnen.

Flensburg, Friesische Str., 86 II. K. Knippenberg.
